% Part: incompleteness
% Chapter: theories-computability
% Section: intepretability

\documentclass[../../../include/open-logic-section]{subfiles}

\begin{document}

\olfileid{inc}{tcp}{itp}

\olsection{$\Th{Q}$ च्या अर्थनिर्धारणाशी सुसंगत उपपत्ती अनिर्णेय असतात}

हे निकाल आणखी सबल करता येतात. अनौपचारिकपणे सांगायचे तर,
भाषा $\Lang{L_1}$ चे दुसरी भाषा $\Lang{L_2}$ मध्ये
अर्थनिर्धारण करताना $\Lang{L_1}$ चे विश्व, संबंधचिन्हे आणि
फलनचिन्हे ही $\Lang{L_2}$ मधील !!{formula}s वापरून
परिभाषित करावी लागतात. येथे त्यासाठी वेळ देणार नसलो,
तरी या संकल्पनेची नेमकी व्याख्या देता येते.

\begin{thm}
  $\Th{T}$ ही अशा भाषेतील उपपत्ती असू द्या की त्या भाषेत
  अंकगणिताच्या भाषेचे अर्थनिर्धारण करता येते, आणि त्या
  अर्थनिर्धारणातील $\Th{Q}$ शी $\Th{T}$ सुसंगत असते.
  मग $\Th{T}$ अनिर्णेय असते. $\Th{T}$ कडून $\Th{Q}$ च्या
  स्वयंसिद्धकांची अर्थनिर्धारित रूपे सिद्ध होत असतील, तर
  $\Th{T}$ चा कोणताही सुसंगत विस्तार निर्णेय नसतो.
\end{thm}

या प्रमेयाची सिद्धता मागील प्रमेयाच्या सिद्धतेत थोडा बदल करून
मिळते; प्रतिउदाहरणावरून $\Th{Q}$ आणि $\Th{\bar Q}$ यांचे
विलगीकरण करता येईल. निवडीच्या स्वयंसिद्धकासह
झर्मेलो--फ्रेंकेल संच उपपत्ती $\Th{ZFC}$ ही सामान्य
गणिताचा मोठा भाग विकसित करण्याइतकी समर्थ स्वयंसिद्धकीय
पायाभूत व्यवस्था मानता येते. $\Th{ZFC}$ मध्ये नैसर्गिक
संख्या परिभाषित करता येतात; आणि या अर्थनिर्धारणातून
$\Th{Q}$ ची स्वयंसिद्धके सत्य ठरतात. म्हणून पुढील निष्कर्ष मिळतो:

\begin{cor}
$\Th{ZFC}$ चा कोणताही सुसंगत, निर्णेय विस्तार नाही.
\end{cor}

\begin{cor}
$\Th{ZFC}$ चा संपूर्ण, सुसंगत आणि संगणनक्षमपणे
!!{axiomatizable} असा कोणताही विस्तार नाही.
\end{cor}

$\Th{ZFC}$ च्या भाषेत $\in$ हे एकमेव द्विस्थानी संबंधचिन्ह
आहे. (खरे तर समतेचीही गरज नाही.) म्हणून पुढील निष्कर्ष मिळतो:

\begin{cor}
द्विस्थानी संबंधचिन्ह असलेल्या कोणत्याही भाषेसाठी प्रथम-क्रम
तर्कशास्त्र अनिर्णेय आहे.
\end{cor}

दोन एकस्थानी फलनचिन्हे असलेल्या कोणत्याही भाषेलाही हा निकाल
लागू होतो, कारण त्यांचा वापर करून द्विस्थानी संबंधचिन्हाचे
अनुकरण करता येते. हे निकाल नेमक्या सीमेवरचे आहेत: केवळ
\emph{एकस्थानी} संबंधचिन्हे आणि जास्तीत जास्त एक
\emph{एकस्थानी} फलनचिन्ह असलेल्या भाषेसाठी प्रथम-क्रम
तर्कशास्त्र निर्णेय असते, असे सिद्ध होते.

आणखी एक माहितीची बाब. प्रमाण प्रतिमानात सत्य असलेल्या
$\Obj 0$, $'$, $+$, $\times$, $<$ या भाषेतील वाक्यांचा
संच अनिर्णेय आहे, हे आपल्याला माहीत आहे. प्रत्यक्षात,
उरलेल्या चिन्हांच्या साहाय्याने $<$ परिभाषित करता येते,
आणि $\times$ व $'$ यांच्या साहाय्याने $+$ परिभाषित
करता येते. म्हणून $\Obj 0$, $'$, $\times$ या भाषेतील
सत्य वाक्यांचा संच अनिर्णेय आहे. दुसरीकडे, प्रमाण
अंकगणितीय प्रतिमानात सत्य असलेल्या $\Obj 0$, $'$, $+$
या भाषेतील वाक्यांचा संच निर्णेय आहे, हे प्रेस्बुर्गर यांनी
दाखवले. मात्र त्याची निर्णयपद्धती संगणनदृष्ट्या अव्यवहार्य आहे.

\end{document}
